\section{Discussion}
\label{sec:discussion}

An archival search can only be as good as the observations it inherits, so
we ask three questions of the result: what new parameter space
the archive allowed us to search, what the absence of a detection actually
constrains, and what stops the search from saying more. The answer to the
last one also says what a dedicated experiment should do differently.

\subsection{What was searched}
\label{sec:newspace}

The parameter space opened here is millimetre and submillimetre frequency
toward the nearest stars. The archive gives \CvUnionA\,GHz of sky frequency
between \SurvFreqLoA{} and \SurvFreqHiA\,GHz toward systems within 40\,pc.
The two \SbPriorFreqGHz\,GHz searches of nearby stars each looked at
a single line. The only search that covered more, that of
\citet{Mason2025}, spanned \MasonBandMHz\,MHz of Band~3, and its targets
were distant stars that happened to lie in calibrator fields
(\S\ref{sec:intro}). Most of this range had therefore not been
searched at all for narrowband emission from the direction of a nearby
star.

The price of using an archive is that the channels were chosen by other
observers, for other purposes. They range from \SzChanFinestKHz\,kHz to
\SzChanWidestMHz\,MHz, while the centimetre-band carrier searches work with
channels of about \SzChanLitHz\,Hz (Fig.~\ref{fig:context}b). A carrier much
narrower than one of these channels is measured against the noise of the
whole channel, and this is one reason why the limits reached here lie well
above those of the centimetre surveys.

Union bandwidth also overstates what was searched, because different
frequencies were observed toward different stars, and to very different
depths. Figure~\ref{fig:waterfall}c therefore counts, at each sky frequency
and EIRP, the systems in which a carrier would have reached the threshold
nine times in ten in a single epoch. At $\CovEirpRefTex$\,W the covered frequency falls to
\CovUnionAtRef\,GHz, \CovFracAtRef{} per cent of the union, and
\CovNSysAtRef{} systems are reached at some frequency, \CovNMaxAtRef{} of them
near \CovNuMaxAtRef\,GHz. However high the power, no frequency is covered
toward more than \CovNMaxFull{} of the \NSysClassA{} systems, and only above
$\CovEirpHalf$\,W is half of the sample reached at any single frequency. The search is thus deep toward a few
systems at a few frequencies, and shallow or absent almost everywhere else.

The \NCoarse{} coarse-channel windows add a constraint of a weaker kind.
Their channels are too wide to follow a drift, so they limit only an
unresolved narrowband spectral excess, but they are numerous, and stacked in
the star's own rest frame they can be deep. Toward Proxima Centauri, the
nearest star, the stacked limit reaches $\SzProxEirpCensus$\,W for an
emitter that stays within one channel in the stellar frame
(\S\ref{sec:stack}).

\subsection{What the non-detection constrains}
\label{sec:excludes}

Every carrier limit in this paper applies to one kind of signal: an unresolved
carrier at the stellar position, confined to a single
\ChanLoA--\ChanHiA\,kHz channel, lying outside the molecular-line mask,
radiating while an archival block was on sky, and drifting linearly within
its window's ceiling of \DriftCeilLo--\DriftCeilHi\,nHz, a fractional drift
rate that corresponds to a line-of-sight acceleration of
\AccelCeilMed--\AccelCeilMax\,m\,s$^{-2}$. That ceiling covers
the orbital acceleration of every known planet in the sample except
\AccelWorstName, which is reached only as precisely as its host's mass and
orbit are known (\S\ref{sec:statistic}). A transmitter on a rapidly rotating
planet, however, uses most of the ceiling on the planet's spin. The Class~B
limits are of a different kind again, on an unresolved spectral excess with
an injection calibration of their own (\S\ref{sec:stack}).

A signal that departs from this model is constrained less, or not at all.
Emission as wide as the widest line in these fields needs \ScNBroadRatio{}
times the power of a one-channel carrier to be found, and a carrier that is
switched on for only part of a block is recovered at that fraction of its
amplitude. Nothing is constrained inside the $\pm\MaskHalfKms$\,km\,s$^{-1}$
mask around each listed transition. Nor does the survey say anything about
the stars within 40\,pc that ALMA has not observed, which are almost all of
the \NGaiaCensus{} that this volume holds.

There are also two blind spots of a more specific kind. The analysis forms
$(\mathrm{XX}+\mathrm{YY})/2$, which retains the full intensity of a
circularly polarised carrier, so what is lost is not sensitivity but the
discrimination that polarisation would offer against unpolarised
astrophysical emission. ALMA's feeds are linear, and the circular signature
lies in the cross-hand products, which this analysis does not extract. The
other blind spot is historical. The \SbPriorFreqGHz\,GHz positronium
hyperfine line that \citet{Steffes1994} and \citet{Mauersberger1996} looked
for lies in ALMA Band~5, where this survey holds \LnBFiveWin{} coarse-channel
windows, all toward \LnBFiveStars, and none of them covers it.

\subsection{What prevents stronger conclusions}
\label{sec:mmdifferent}

Depth is not the main limit. As an illustration, a \BenchDiam\,m dish
radiating \BenchPowerMW\,MW at \BenchFreqGHz\,GHz with an aperture efficiency
of \BenchEta{} has a directive gain $G=\eta(\pi D\nu/c)^{2}=\BenchGain$, and
hence an EIRP of $\BenchEirp$\,W. The gain grows as the square of the
frequency, and the systems' best windows lie between \BenchTunedFreqLo{} and
\BenchTunedFreqHi\,GHz, so the fair comparison is with the same dish tuned to
each system's own best window, where its EIRP runs from
$\BenchTunedEirpLo$ to $\BenchTunedEirpHi$\,W. Tuned in that way it lies
above the single-epoch limit of \BenchNSysTuned{} of the \NSysALimSum{}
Class~A systems that retain unmasked bandwidth and have a bounded position,
so a transmitter of terrestrial scale would already have been within reach
toward those systems. For comparison, the Arecibo planetary radar had an EIRP
of about $\AroEirp$\,W at \AroFreqGHz\,GHz \citep{EarthDetectingEarth2025},
about \AroMedMult{} times below the median limit, and \AroNSysTen{}
systems come within a factor of ten of it. Near that boundary
\AroTenEdgeClause. The dish's beam, however, has a solid
angle $\Omega=4\pi/G=\BmOmegaSr$\,sr and covers only $\BmSkyFrac$ of the sky.
EIRP is not transmitter power, and these limits apply to a beamed
transmitter only if its beam happens to contain the Solar System. The
archive is limited less by its depth than by the astrophysics that can
imitate a carrier, and by the repeat coverage that it happens to hold.

At these wavelengths the false positive is line emission. A
catalogued molecular transition lies within the mask for \CsNAttrQp{} of the
\CsNCrossQp{} crossings in the census's quality-passing data. For
\LnAttrSupp{} of them there is support beyond the frequency match: a second
transition of the same species, catalogued gas on the sightline, or the
crossing's own recurrence. The other \LnAttrCoincWord{} rest on the
coincidence alone, which is not an identification. None was identified as
interference by the available diagnostics (Appendix~\ref{app:rfi}), which
leaves \CsNUnattrQp{} for the comparison with chance of \S\ref{sec:trigunif}.

Molecular lines also bring a test of their own that has no centimetre
analogue. A line that fills the primary beam raises the star and its control
annulus together, whereas a carrier from the star raises only the star. Of the \LrStageN{}
crossings whose statistic at the star exceeds that at every control position
of their window, which is the test that singles out the field leader
(\S\ref{sec:vistest}), \LrStageNCo{} are carbon monoxide toward
\LrStageCoStars{} (\S\ref{sec:bpicval}). All of them are listed in
\texttt{\LrStageFile}.

An archive supplies repeat epochs only by accident. A single crossing,
however strong, becomes a testable claim only when the star is observed
again in a separated session, at the same frequency for an emitter at rest in
its frame or at the frequency a trajectory predicts for one that drifts, and
the archive holds such repeats only where observers happened to return.
We describe this by the same-frequency coverage
$C_{\rm conf}(\nu)$, the number of systems toward which two separated blocks
(Table~\ref{tab:quantities}) both cover the sky frequency $\nu$. It measures
coverage, and is not a completeness for an orbiting carrier, because a
drifting carrier must also fall at its predicted frequency in the second
block. Its support is \CvConfGHzDay\,GHz, \CvConfPctDay{} per cent of the
\CvUnionA\,GHz searched, and only \CvConfGHzYr\,GHz beyond a year
(Fig.~\ref{fig:waterfall}b). Of the \NSysClassA{} systems, \ZcConfSysDay{}
have such a block and \ZcConfSysYr{} have one more than a year away. The
other \ZcNoRepSys{} have none, for the two reasons set out in
\S\ref{sec:conclusions}. Even where a repeat exists it can be
too shallow to test anything, as it was for \HjNPowerLo{} of the \RcNExcl{}
crossings read at a predicted cell.

\subsection{What a dedicated ALMA experiment should do differently}
\label{sec:future}

Each of these limits has a remedy on the same telescope. A dedicated search
can choose its targets and its frequencies instead of inheriting them. This
archive favours dusty discs, which are the worst case for line confusion, and
covers least the M~dwarfs that host most of the nearby habitable-zone
planets. Fine channels are needed in every window, since the \NCoarse{}
coarse-channel windows of the archive cannot carry a carrier search at all.
Tunings placed away from the bright molecular lines recover most of the
\MkMaskAGHz\,GHz of Class~A coverage that the mask costs here, and probably
remove many of the false positives as well. Above all, the recurrence
observations have to be scheduled deliberately, at the same tuning and days
to years apart, so that every crossing meets a separated epoch at the
frequency it predicts. This is chiefly a matter of scheduling, and each
additional epoch tests a trajectory predicted independently from the others.
Indeed \PaNSysThree{} of the \NSysClassA{} systems already carry three epochs
at a common frequency, but none of them has a crossing seen in more than one
of those epochs. The obvious first target for such a programme is \BkAlmaName, the most
discrepant of the objects whose repeat windows were compared with their own
control positions (\S\ref{sec:twostars}).
